% Work-MAT | 2. rocnik | pripravovana hodina 43
% Návaznost: uživatelské připravy.zip (hodiny 01–20), autorská učebnice na GitHubu.
% Realisticky.cz: didaktická inspirace v posloupnosti a typech úloh; zadání jsou vlastní.
\documentclass[11pt,a4paper]{article}
\input{../../sablony/pripravy/priprava.sty}
\nastavPripravu{2. ročník}{Logaritmické nerovnice}{43}{Logaritmické nerovnice II}
\begin{document}
\blok{Motivace a návaznost}
Porovnej rostoucí a klesající logaritmus: změna základu neznamená převrácení funkčních hodnot, ale monotónnosti. Počítej i otevřené konce intervalů.
\poznamka{Orientační nabídka pro 45--55 minut: krátké zopakování a motivace, nový výklad, tři společné vzorové úlohy, samostatná práce. Vyber podle tempa třídy; zbývající zadání slouží jako rezerva.}
\blok{Výklad}
Pro $0<a<1$ platí $\log_a u>\log_a v\iff u<v$ při $u,v>0$.
\blok{Společné příklady}
\begin{prikladReseniBox}[title={Příklad 1}]
Vyřeš, resp. urči: $\log_{\frac{1}{2}}(x+2)\ge-2$.
\tcblower
\[\begin{aligned}x>-2 \\ x+2\le4 \\ -2<x\le2\end{aligned}\]
\end{prikladReseniBox}
\begin{prikladReseniBox}[title={Příklad 2}]
Vyřeš, resp. urči: $\log_{\frac{1}{3}}(x-1)<-1$.
\tcblower
\[\begin{aligned}x>1 \\ x-1>3 \\ x>4\end{aligned}\]
\end{prikladReseniBox}
\poznamka{V samotném přechodu na argumenty u klesajícího logaritmu obracej směr. Různá znaménka jsou hlavní riziko.}
\blok{Samostatné procvičování}
\begin{zadaniReseniBox}[title={Samostatná práce 1}]
\begin{enumerate}[label=\alph*)]
\item $\log_{\frac{1}{2}}(x+1)\ge -1$.
\item $\log_{\frac{1}{2}}(x+2)\ge -2$.
\item $\log_{\frac{1}{2}}(x+3)\ge -3$.
\item $\log_{\frac{1}{2}}(x+4)\ge -1$.
\item $\log_{\frac{1}{2}}(x+1)\ge -2$.
\item $\log_{\frac{1}{2}}(x+2)\ge -3$.
\end{enumerate}
\tcblower
\begin{enumerate}[label=\alph*)]
\item $-1<x\le1$.
\item $-2<x\le2$.
\item $-3<x\le5$.
\item $-4<x\le-2$.
\item $-1<x\le3$.
\item $-2<x\le6$.
\end{enumerate}
\end{zadaniReseniBox}
\blok{Dobrovolné rozšíření}
\begin{zadaniReseniBox}[title={Rozšíření 1}]
\begin{enumerate}[label=\alph*)]
\item $Vytvoř vlastní úlohu tohoto typu a se spolužákem porovnej podmínky a výsledek.$.
\end{enumerate}
\tcblower
\begin{enumerate}[label=\alph*)]
\item Zkontroluj definice a argumenty..
\end{enumerate}
\end{zadaniReseniBox}
\blok{Úkol na doma}
U každého logaritmu vypiš podmínky argumentu a proveď zkoušku. U rovnic vyznač případně nepřípustné kořeny, u nerovnic odpověz intervalem.
\begin{zadaniReseniBox}[title={Domácí úkol}]
\begin{enumerate}[label=\alph*)]
\item $\log_{\frac{1}{2}}(x+3)\ge -1$.
\item $\log_{\frac{1}{2}}(x+4)\ge -2$.
\item $\log_{\frac{1}{2}}(x+1)\ge -3$.
\item $\log_{\frac{1}{2}}(x+2)\ge -1$.
\item $\log_{\frac{1}{2}}(x+3)\ge -2$.
\item $\log_{\frac{1}{2}}(x+4)\ge -3$.
\item $\log_{\frac{1}{2}}(x+1)\ge -1$.
\item $\log_{\frac{1}{2}}(x+2)\ge -2$.
\end{enumerate}
\tcblower
\begin{enumerate}[label=\alph*)]
\item $-3<x\le-1$.
\item $-4<x\le0$.
\item $-1<x\le7$.
\item $-2<x\le0$.
\item $-3<x\le1$.
\item $-4<x\le4$.
\item $-1<x\le1$.
\item $-2<x\le2$.
\end{enumerate}
\end{zadaniReseniBox}
\poznamka{Domácí úkol obsahuje 8 položek; předpokládaná domácí práce přibližně 15--30 minut. Po větších blocích jsou vědomě promíchané dříve probrané typy.}
\end{document}
